\documentclass[10pt,a4paper]{amsart}
\usepackage[margin=19mm]{geometry}
\usepackage{amsmath,amssymb,mathtools}
\usepackage[scr=rsfs,cal=euler]{mathalfa}
\usepackage{xfrac}
\usepackage[unicode,hidelinks,bookmarks=false]{hyperref}
\newcommand{\ie}{i.e.\ }
\newcommand{\cf}{cf.\ }
\newcommand{\resp}{respectively\ }
\newcommand{\Subsection}{\subsection}
\providecommand{\nobreakdash}{\nobreak\hbox{-}}
\setlength{\emergencystretch}{2em}
\multlinegap0pt
\usepackage{amssymb}
\usepackage{enumerate}
\usepackage{xcolor}
\newtheorem{lemma}{Lemma}
\newtheorem{itemizedot}{Itemizedot}
\newtheorem{enumeratenumeric}{Enumeratenumeric}
\title{Nonuniqueness of Leray--Hopf solutions to the unforced incompressible 3D Navier--Stokes Equation}
\author{Thomas Hou, Yixuan Wang, Changhe Yang}
\date{Source: arXiv:2509.25116v2 (CC BY 4.0)}
\begin{document}
\maketitle

\noindent\emph{Source and licence.} Nonuniqueness of Leray--Hopf solutions to the unforced incompressible 3D Navier--Stokes Equation. Version arXiv:2509.25116v2 (CC BY 4.0), distributed by arXiv under the Creative Commons Attribution 4.0 International licence, http://creativecommons.org/licenses/by/4.0/, as recorded in the arXiv licence metadata for this exact version and shown on the abstract page https://arxiv.org/abs/2509.25116v2. Full licence notice: \texttt{licenses/arxiv-2509.25116-CC-BY-4.0.txt}.\par\smallskip

\noindent\textit{Editorial scope.} Counted unit: the article's lemma lem:zeros (source0.tex lines 3394--3425). The article's complete original proof of it is delivered with it (source0.tex lines 4308--4386, one \texttt{proof} environment of 79 lines, which the article places away from the statement). No mathematics is written, repaired or re-derived: the statement and the proof are the article's own lines, in the article's own notation, and no retained line is substituted or re-flowed.  Retained uncounted as the source-local context the proof needs: lines 3325--3328; lines 3336--3339.  Nothing was omitted from any retained span, so the package carries no omission record.  No retained line contains a citation, so the wrapper carries no bibliography and no bibliography entry.  No retained line contains a figure environment, an image-inclusion command or drawing code, so no image is delivered and the package declares no asset dependency.  Editorial formatting only, in addition: an AMS \texttt{amsart} preamble that restates the article's own declarations and macros used by the retained lines, namely the environments \texttt{lemma}, \texttt{itemizedot}, \texttt{enumeratenumeric} on separate counters, together with the standard packages those lines need.  The retained environments print in this document as Lemma 1, Itemizedot 1, Enumeratenumeric 1 in order of appearance, whereas the article prints the same items with the numbering of its own full document.  The complete native source of the article as distributed in the arXiv e-print for this exact version is bundled unchanged as \texttt{sources/186000/source0.tex}.
  \begin{equation}
    rJ_{l + \frac{1}{2}}' (r) - \frac{1}{2} J_{l + \frac{1}{2}} (r) = 0,
    \label{eq:alphaZ}
  \end{equation}

  \begin{equation}
    (r^2 - l (l - 1) (l + 2)) r \tilde{j}_l' (r) + (r^2 - l (l - 1))  (r^2 -
    (l - 1) (l + 2))  \tilde{j}_l (r) = 0, \label{eq:alphaY}
  \end{equation}

\begin{lemma}
  \label{lem:zeros}Let $0 = \beta_{l, 0} < \beta_{l, 1} < \beta_{l, 2} <
  \cdots$ denote the zeros of $J_{l + \frac{1}{2}} (k)$ and $\gamma = \sqrt{l
  (l - 1)  (l + 2)}$.
  \begin{itemizedot}
    \item If $l = 1$, define $\tilde{\beta}_{1, j} = \beta_{1, j}$;
    
    \item If $l \geqslant 2$ and $\gamma \neq \beta_{l, j}$ for any $j > 0$,
    denote $\tilde{\beta}_{l, j}$ as the sequence obtained by sorting
    $\{\beta_{l, j}, \gamma\}$ in an ascending order.
  \end{itemizedot}
  Let $s_l$ be the unique positive root of:
  \[ Q_l (s) = s^3 - 4 l^2 s^2 + l (l - 1)  (6 l^2 + 11 l + 7) s - 3 l (l -
     1)^2  (l + 1)^2  (l + 2), \]
  and $t_l = \sqrt{l (l + 1)}$. Assume that $\sqrt{s_l}$ is not a root of
  \eqref{eq:alphaY} and that $t_l$ is not a root of \eqref{eq:alphaZ}.
  
  Then the following statements hold:
  \begin{enumeratenumeric}
    \item The equation \eqref{eq:alphaY} has exactly one root on every
    interval $(\tilde{\beta}_{l, j}, \tilde{\beta}_{l, j + 1})$ for $j
    \geqslant 0$.
    
    \item The equation \eqref{eq:alphaZ} has exactly one root on every
    interval $(\beta_{l, j}, \beta_{l, j + 1})$ for $j \geqslant 0$.
    
    \item The first zeros $\alpha_{l, 1}^Y$ and $\alpha_{l, 1}^Z$ satisfy the
    lower bounds:
    \[ \alpha_{l, 1}^Y \geqslant \sqrt{s_l}, \quad \alpha_{l, 1}^Z \geqslant
       t_l . \]
  \end{enumeratenumeric}
\end{lemma}

\begin{proof}[Proof of Lemma \ref{lem:zeros}]
  First we consider the equation \eqref{eq:alphaY}. Since $J_{l + \frac{1}{2}}
  (\gamma) \neq 0$ for $l \geqslant 2$, it can be written in the form:
  \[ \frac{kJ_{l + \frac{1}{2}}' (k)}{J_{l + \frac{1}{2}} (k)} - \frac{3}{2} +
     \frac{(k^2 - l (l - 1))  (k^2 - (l - 1) (l + 2))}{(k^2 - l (l - 1) (l +
     2))} = 0. \]
  If $l = 1$, the equation becomes:
  \[ \frac{kJ_{l + \frac{1}{2}}' (k)}{J_{l + \frac{1}{2}} (k)} - \frac{3}{2} +
     k^2 = 0. \]
  Define:
  \[ R_l (k) = \frac{kJ_{l + \frac{1}{2}}' (k)}{J_{l + \frac{1}{2}} (k)},
     \quad T_l (k) = \frac{(k^2 - l (l - 1))  (k^2 - (l - 1) (l + 2))}{(k^2 -
     l (l - 1) (l + 2))}, \quad F_l (k) = R_l (k) - \frac{3}{2} + T_l (k) . \]
  We observe that $F_l (k)$ is continuous except at $\tilde{\beta}_{l, j}$ ($j
  \geqslant 1$). If $\tilde{\beta}_{l, j}$ is a zero of $J_{l + \frac{1}{2}}
  (x)$, then
  \[ \lim_{x \rightarrow \tilde{\beta}_{l, j} +} R_l (x) = + \infty, \quad
     \lim_{x \rightarrow \tilde{\beta}_{l, j} -} R_l (x) = - \infty, \]
  and $T_l (\tilde{\beta}_{l, j})$ is finite. If $\tilde{\beta}_{l, j}$ is
  $\gamma$, we have
  \[ \lim_{x \rightarrow \tilde{\beta}_{l, j} +} T_l (x) = + \infty, \quad
     \lim_{x \rightarrow \tilde{\beta}_{l, j} -} T_l (x) = - \infty, \]
  and $R_l (\tilde{\beta}_{l, j})$ is finite. Therefore, for every
  $\tilde{\beta}_{l, j}$, in both cases we obtain 
  \[ \lim_{x \rightarrow \tilde{\beta}_{l, j} +} F_l (x) = + \infty, \quad
     \lim_{x \rightarrow \tilde{\beta}_{l, j} -} F_l (x) = - \infty . \]
  At $\tilde{\beta}_{l, 0} = 0$, the power series
  expansion
  \[ J_{l + \frac{1}{2}} (x) = \frac{(x / 2)^{l + \frac{1}{2}}}{\Gamma (l +
     \frac{3}{2})}  (1 - \frac{x^2}{4 (l + \frac{3}{2})} + O (x^4)) \]
  shows that
  \[ F_l (k) = \frac{3 (1 + l)^2}{l (2 + l)  (3 + 2 l)} k^2 + o (k^2) . \]
  Therefore, we obtain $F_l (k) > 0$ for sufficiently small $k$. We conclude that $F_l
  (k)$ has at least one root on every interval $(\tilde{\beta}_{l, j},
  \tilde{\beta}_{l, j + 1})$ with $j \geqslant 0$. The case $l = 1$ follows by
  a similar argument.
  
  Next we derive a lower bound for the first root, and prove that assuming two
  zeros in the same interval leads to a contradiction. Since $F_l 
  (\tilde{\beta}_{l, j} +) = \infty$ and $F_l  (\tilde{\beta}_{l + 1, j} -) =
  - \infty$, we denote $k_2 > k_1$ as the first two roots in
  $(\tilde{\beta}_{l, j}, \tilde{\beta}_{l, j + 1})$, which implies that $F_l'
  (k_1) \leqslant 0$.
  
  Straightforward calculation shows that
  \[ R_l' (k) = - \frac{R_l (k)^2 + k^2 - \left( l + \frac{1}{2}
     \right)^2}{k}, \quad T'_l (k) = 2 k \left( 1 - \frac{l (l - 1)^3 (l + 1)
     (l + 2)}{(k^2 - l (l - 1) (l + 2))^2} \right) . \]
  When $F_l (k) = 0$, we have $R_l (k) = \frac{3}{2} - T_l (k)$. Thus $F_l'$
  simplifies into:
  \[ F_l' (k) = - \frac{\left( \frac{3}{2} - T_l (k) \right)^2 + k^2 - \left(
     l + \frac{1}{2} \right)^2}{k} + T'_l (k) = - \frac{kQ_l (k^2)}{(k^2 - l
     (l - 1) (l + 2))^2}, \]
  where
  \[ Q_l (s) = s^3 - 4 l^2 s^2 + l (l - 1)  (6 l^2 + 11 l + 7) s - 3 l (l -
     1)^2  (l + 1)^2  (l + 2) . \]
  It is easy to verify that for $l \geqslant 2$, we have by quadratic
  equations that
  \[ Q_l' (s) = 3 s^2 - 8 l^2 s + l (l - 1)  (3 l + 7)  (2 l + 1) > 0. \]
  Therefore, $Q_l (s)$ has a unique root $s_l > 0$ such that $Q_l (s) < 0$ for
  $s < s_l$, and $Q_l (s)$ is positive and strictly increasing for $s > s_l$.
  
  For $l = 1$, we have $Q_1 (s) = s^3 - 4 s^2$, which satisfies the same
  property.
  
  Since $F_l' (k_1) \leqslant 0$, it follows that
  \[ k_1 \geqslant \sqrt{s_l}, \]
  and that $0 \leqslant Q_l (k_1^2) < Q_l (k_2^2)$. Hence, we obtain $F_l'
  (k_2) < 0$. Thus $F_l$ cannot cross the axis at $k_1$, and it has to hold
  that $F_l (k_1) = F_l' (k_1) = 0$. We verify that $F_l (\sqrt{s_l}) \neq 0$.
  Hence, we reach a contradiction and conclude the proof of the lemma.
  
  By a similar but simpler argument, the equation \eqref{eq:alphaZ} has
  exactly one root in every interval $(\beta_{l, j}, \beta_{l, j + 1})$ for
  all $l \geqslant 1, j \geqslant 0$. Moreover the first root satisfies the
  lower bound
  \[ \beta_{l, 1} \geqslant \sqrt{l (l + 1)}, \]
  provided that $\sqrt{l (l + 1)}$ is not a root.
\end{proof}

\end{document}
